% =========================================================
% Structural Induction (Preview)
% =========================================================
\paragraph{Structural Induction (Preview).}

Induction on $\mathbb{N}$ is a special case of a general principle
that applies to any recursively defined structure. You do not need
$\mathbb{N}$ to do induction; you need a recursive definition.

\begin{remark}[The general principle]
Let $X$ be a set defined by:
\begin{enumerate}[itemsep=2pt]
  \item A collection of \emph{base cases} (atomic elements not built
    from other elements), and
  \item A collection of \emph{constructor rules} (ways of building
    larger elements from smaller ones).
\end{enumerate}
Structural induction on $X$: to prove $P(x)$ for all $x \in X$,
prove $P$ for all base cases, then prove that if $P$ holds for the
parts, it holds for the whole.

This is exactly standard induction where the ``less than'' ordering
on $\mathbb{N}$ is replaced by the ``proper sub-element'' ordering
on $X$.
\end{remark}

\begin{example}[Induction on algebraic terms]
Define an \emph{algebraic term} over a ring $R$ by:
\begin{itemize}[itemsep=2pt]
  \item Base case: any element $r \in R$ is a term.
  \item Constructor: if $s, t$ are terms, so are $s + t$ and $s \cdot t$.
\end{itemize}
To prove a property $P$ holds for all terms:
\begin{enumerate}[itemsep=2pt]
  \item Prove $P(r)$ for all $r \in R$ (base cases).
  \item Prove: if $P(s)$ and $P(t)$ hold, then $P(s + t)$ holds.
  \item Prove: if $P(s)$ and $P(t)$ hold, then $P(s \cdot t)$ holds.
\end{enumerate}
The recursive structure of terms determines the structure of the proof.
\end{example}

\begin{remark*}[Exposition]
Structural induction is the proof principle paired with structural recursion.
Both follow the constructors of the object being studied. Structural recursion
defines a function by giving its value on the base cases and then giving its
value on each constructed object in terms of the values already defined on the
immediate parts. Structural induction proves a predicate or property by showing
that it holds for the base cases and is preserved by every constructor.

Thus structural induction is like ordinary mathematical induction: it proves a
statement $P(x)$ for every object $x$ in a recursively generated class. The
difference is that the next step is not always the Peano successor
$n\mapsto n\pp$. The next step is whatever constructor builds the object. For
natural numbers the constructors are $0$ and successor. For formulas the
constructors are atomic formulas, negation, binary connectives, and
quantifiers. For trees the constructors are leaves and nodes.

A useful way to see this is to attach a syntax tree to each well-formed
formula. Define trees for propositional formulas recursively:
\begin{itemize}[itemsep=2pt]
  \item if $P\in\Prop$, then $\operatorname{leaf}(P)$ is a tree;
  \item if $T$ is the tree for $\varphi$, then
    $\operatorname{node}(\neg,T)$ is the tree for $\neg\varphi$;
  \item if $T_\varphi$ and $T_\psi$ are the trees for $\varphi$ and $\psi$,
    then $\operatorname{node}(\circ,T_\varphi,T_\psi)$ is the tree for
    $(\varphi\circ\psi)$.
\end{itemize}
Operations on this tree are recursive in the same sense. The depth of a leaf is
$0$; the depth of a negation node is one plus the depth of its child; the depth
of a binary node is one plus the maximum of the depths of its two children.
Structural induction then proves statements about all such trees, and hence
about all formulas, by checking the leaf case, the negation-node case, and the
binary-node case. The proof follows the same recursive blueprint as the
definition.
\end{remark*}

\begin{remark}[Where structural induction appears in this project]
Induction on $\mathbb{N}$ is structural induction on the Peano
structure $(\mathbb{N}, 0, S)$: the base case is $n = 0$, the
constructor is $n \mapsto n\pp$.

In the Abstract Algebra chapter, polynomial properties are proved by
induction on degree, which is structural induction on the recursive
definition of polynomials over a ring.

In computer science, correctness of recursive algorithms is proved by
structural induction on the input data structure (lists, trees,
expressions). The full treatment of structural induction appears in
the Abstract Algebra chapter.
\end{remark}
